В ходе выполнения данной магистерской выпускной квалификационной работы успешно разработана и верифицирована гибридная система для авто\-ма\-ти\-зи\-ро\-ван\-но\-го разбора библиографических ссылок из научных публикаций. Актуальность иссле\-до\-ва\-ния обусловлена острой необходимостью перехода от трудоемкой ручной обработки текстов к автоматическим методам в современных сис\-те\-мах анализа научной дея\-тель\-нос\-ти (таких как ИСАНД), а также специфическими трудностями извлечения структурированных данных. К ним относятся деструктивные факторы разметки: ар\-те\-фак\-ты конвертации PDF-слоя, ошибки OCR и \guillemotleft стилевой хаос\guillemotright{} оформления источников.

Задача извлечения библиографии в работе декомпозирована на три последовательных этапа: локализация области библиографии, построчная сегментация записей и се\-ман\-ти\-чес\-кое структурирование атомарных полей. Для оценки качества на всех этапах конвейера обоснованно выбрана символьно-ориентированная метрика $F_1$. Данный подход позволяет корректно учитывать частичные перекрытия границ текстовых сущностей, возникающие из-за сдвига координат и отсутствия пробелов при парсинге PDF-документов.

Анализ результатов численных экспериментов на открытом англоязычном наборе данных и специализированном русскоязычном корпусе показал, что универсального метода, оптимального для всех этапов пайплайна, не существует. Для макро-ло\-ка\-ли\-за\-ции области библиографии наилучшие результаты продемонстрировали регулярные выражения. На этапе извлечения сырых библиографических записей наиболее эф\-фек\-тив\-ным среди рассмотренных методов оказался инструмент GROBID, достигший значения $F_1 = 0{,}97$. С задачей семантического структурирования полей из отдельных записей наилучшим образом справилась большая языковая модель Llama-3-8B без использования LoRA-ад\-ап\-те\-ров, продемонстрировав значение макро-$F_1 = 0{,}84$.

На основе полученных сравнительных характеристик спроектирована и программно реализована гибридная архитектура. В ней инструмент GROBID отвечает за сегментацию документа на отдельные записи, а модель Llama-3-8B осуществляет их финальное извлечение в структурированный формат JSON. Разработанный программный модуль успешно прототипирован в состав ИСАНД для автоматического построения кас\-то\-ми\-зи\-ро\-ван\-ных библиометрических сетей. Результаты работы прошли апробацию на профильных конференциях, а теоретические положения опубликованы в журнале \guillemotleft Системы высокой доступности\guillemotright~\cite{system_higher}.

Для дальнейших исследований остается открытым ряд научно-практических задач. Используемый на этапе сегментации инструмент GROBID ориентирован на английский язык, что требует его глубокого дообучения на русскоязычном корпусе или замены на альтернативные отечественные решения. Необходима модернизация системы для парсинга источников, заданных не в финальном списке литературы, а в постраничных и концевых сносках, характерных для гуманитарных наук. Требует решения новая фундаментальная проблема появления \guillemotleft генеративных галлюцинаций\guillemotright{} в списках литературы -- вымышленных авторами с помощью ИИ-ассистентов библиографических записей, верификация которых возможна только путем интеграции модулей внешнего семантического соответствия с API наукометрических баз Crossref и РИНЦ.